\section{Joint identification and the full sequence}\label{joint:section}

The preceding identifications were made on a common contour extraction.
We now show that their joint law is the one in \cref{intro:main}.
The only remaining coordinate choice is a single possible reflection
of the entire tuple; its sign need not be independent of that tuple.

\begin{proof}[Proof of \cref{intro:main}]
Keep the deterministic normalization $a_n$ of
\cref{intro:intrinsic-input}, and start with an arbitrary deterministic
subsequence of positive integers. By
\cref{enc:contour-limits,enc:peanosphere}, pass to a further joint
extraction of the contours and their reference curve-decorated sphere.
Retain the compatible local arrays used in the conformal, loop and
metric arguments. The compactified observation spaces permit this
joint extraction, and their conditional-law identifications remain
valid with the additional observations. Write $h^{\mathrm o}$,
$\mu^{\mathrm o}$ and $D^{\mathrm o}$ for the field, area and metric
in the original oriented reference coordinate, and
$\Gamma^{\mathrm o}$ for its complete nested loop ensemble.

The homeomorphisms $H_n$ of \cref{top:homeomorphism} are chosen before
the three fresh finite area samples. The reference area convergence
$(H_n)_*m_n\Rightarrow\mu^{\mathrm o}$ is established in
\cref{conf:canonical}. By \cref{joint:marks}, the sample images
converge jointly to three points with conditional law
$(\mu^{\mathrm o})^{\otimes3}$ given the reference surface and its
exploration. Let $T$ be the orientation-preserving
M\"obius map sending these three limiting samples to $0,1,\infty$.
\Cref{conf:canonical,conf:mesh-area} give, after further extraction,
\[
 f_n:=\phi_n\circ H_n^{-1}\longrightarrow f
 \quad\hbox{uniformly},\qquad
 f\in\{T,\mathfrak r\circ T\},
 \qquad \mathfrak r(z)=\overline z.
\]
Here $\mathfrak r(\infty)=\infty$. On this same coupling,
$\mu_n:=(\phi_n)_*m_n$ converges weakly to $f_*\mu^{\mathrm o}$,
and the maximum spherical diameter of a flag image tends to zero.

\Cref{loop:matching,loop:canonical} give complete loop matchings on
this extraction. Every loop of prescribed positive diameter on either
side is matched, with uniform error tending to zero up to
reparameterization. Thus $\Gamma_n$ converges to
$f\Gamma^{\mathrm o}$, retaining cyclic order, multiplicity and all
nesting. The map $f$ is exactly the map used for the area limit.

The distance graphs are tight jointly with these data. Their first two
coordinates lie in the compact sphere, and their heights are bounded
by $a_n\diam(V(M_n),d_n)$, which is tight by
\cref{intro:intrinsic-input}. Below any fixed height $L$, the
nonempty compact subsets of $\Chat^2\times[0,L]$ form a compact
Hausdorff hyperspace. Retain a convergent further subsequence of these
graphs and of the intrinsic spaces. \Cref{comp:metric-graph} identifies
its graph limit, sample by sample, as
\[
 \{(f(z),f(w),D^{\mathrm o}(z,w)):z,w\in\Chat\}.
\]
This identification includes all vertex pairs, including pairs
approaching the exploration root. Its scalar is one for the prescribed
$a_n$: the completion argument obtains a deterministic scalar multiple
of the reference metric and fixes that scalar using the intrinsic
diameter law. It makes no change to the retained area or loop coupling.

Consider first the orientation-preserving target tuple
\[
 Y=\left(T_*\mu^{\mathrm o},
   \{(T(z),T(w),D^{\mathrm o}(z,w)):z,w\in\Chat\},
   T\Gamma^{\mathrm o}\right).
\]
The first two entries are the area and intrinsic metric of the ordinary
unit-area sphere in its three-area-mark coordinate. Conditional on the
reference surface, $\Gamma^{\mathrm o}$ has the fixed whole-plane
nested CLE law, by \cref{loop:canonical}. The three marks use fresh
sampling randomness, so for every bounded measurable loop functional
$G$,
\[
 \EE[G(T\Gamma^{\mathrm o})\mid
           \text{reference surface and three marks}]
 =\int G(T\mathcal L)\,\dd\mathsf{CLE}_{\kappa}(\mathcal L)
 =\int G(\mathcal L)\,\dd\mathsf{CLE}_{\kappa}(\mathcal L).
\]
The second equality is M\"obius invariance. Taking conditional
expectation given the normalized embedded marked surface proves the
independence required by \cref{intro:main}. Thus $Y$ has exactly the
asserted target law.

Let $X$ be the limiting discrete tuple on this coupling. The three
identifications above give $X\in\{Y,\mathfrak r Y\}$, with reflection
acting simultaneously on the measure, both spatial coordinates of the
distance graph, and every loop. Both tuple laws are reflection invariant.
For the discrete tuple, reflection is a weight-preserving bijection of
the oriented finite-map ensemble. It preserves flag areas and all-edge
graph distances and transports the independent area samples. Since
$0,1,\infty$ are fixed by conjugation, the normalized tuple is reflected
as well. Reflection is continuous in all three topologies, so this
invariance passes to $X$. The law of $Y$ is invariant by reflection
symmetry of the ordinary quantum sphere, its area and metric, and the
independent unoriented CLE ensemble. \Cref{joint:reflection} now gives
$\Law(X)=\Law(Y)$, even when the sign of $f$ depends on all the retained
data.

Every deterministic subsequence therefore has a further jointly
convergent subsequence with the same specified law. The subsequence
criterion proves \eqref{intro:joint-convergence} through all positive
integers. On each such extraction the maximum flag diameter also tends
to zero. If \eqref{intro:mesh-conclusion} failed in probability, a
subsequence would have a fixed positive probability of exceeding a
fixed positive threshold, contradicting its further extraction from
\cref{conf:mesh-area}. This proves the mesh conclusion and completes
the theorem.
\end{proof}
